\section{Introduction}
\label{sec:introduction}
In this paper, we study the structured nonsmooth nonconvex composite problem
\begin{align}
    \min_{\bx\in\RR^n}\quad & F(\bx) := f(\bx) + h(\calA(\bx)),
    \tag{\textbf{P}}
    \label{prob:main}
\end{align}
where $f\colon\RR^n\to\RR$ is a possibly nonconvex function with a Lipschitz
continuous gradient, $\calA\colon\RR^n\to\RR^m$ is a smooth and possibly nonlinear
mapping, and $h\colon\RR^m\to\RR\cup\{+\infty\}$ is a proper closed convex function
whose proximal mapping is easy to evaluate; the complete assumptions are stated in
Section~\ref{sec:assumptions}. The two components play complementary roles: $f$ is
usually a smooth loss or data-fidelity term, whereas $h$ is a nonsmooth convex function
that imposes robustness, sparsity, or other structure on the transformed variable
$\calA(\bx)$
\cite{charisopoulos2021low,markovsky2008structured,parekh2015convex}. Several models
of practical interest take this form. 

\subsection{Motivating Examples}

\textbf{Robust phase retrieval}~\cite{duchi2019solving}.
Given measurement vectors $\mathbf a_i\in\RR^n$ and observations
$b_i\approx\langle\mathbf a_i,\bx^\star\rangle^2$, a fraction of which may be grossly
corrupted, one solves
\begin{equation*}
  \min_{\bx\in\RR^n}\;\frac1m\sum_{i=1}^m
  \bigl|\langle\mathbf a_i,\bx\rangle^2-b_i\bigr|.
\end{equation*}
This is~Problem \eqref{prob:main} with $f\equiv0$, $h=\frac1m\|\cdot\|_1$, and
$\calA(\bx)_i=\langle\mathbf a_i,\bx\rangle^2-b_i$; the $\ell_1$ loss makes the
model robust to outliers.

\textbf{Robust low-rank matrix recovery}~\cite{charisopoulos2021low}.
Given a linear measurement map $\mathcal M\colon\RR^{d_1\times d_2}\to\RR^m$ and
observations $\mathbf b$ containing outliers, a rank-$r$ matrix is recovered in
factorized form via
\begin{equation*}
  \min_{\mathbf U\in\RR^{d_1\times r},\,\mathbf V\in\RR^{d_2\times r}}
  \frac1m\bigl\|\mathcal M(\mathbf U\mathbf V^\top)-\mathbf b\bigr\|_1.
\end{equation*}
Here $\bx=(\mathbf U,\mathbf V)$, $h=\frac1m\|\cdot\|_1$, and
$\calA(\bx)=\mathcal M(\mathbf U\mathbf V^\top)-\mathbf b$ is bilinear in
$(\mathbf U,\mathbf V)$, so the model again has the form~Problem \eqref{prob:main}.

\textbf{Minimizing a pointwise maximum}~\cite{zhang2020single}.
For smooth losses $f_1,\dots,f_m$, such as the losses on $m$ data subgroups, one solves
\begin{equation*}
  \min_{\bx\in\RR^n}\;\max_{1\le i\le m} f_i(\bx).
\end{equation*}
This is~Problem \eqref{prob:main} with $h(u)=\max_{1\le i\le m}u_i$ and
$\calA(\bx)=(f_1(\bx),\dots,f_m(\bx))$; a smooth regularizer, if present, enters
through $f$, and the proximal mapping of $h$ reduces to a projection onto the unit
simplex.

\subsection{Related Work}

Introducing the auxiliary variable $\by=\calA(\bx)$ rewrites~Problem \eqref{prob:main} as
the equality-constrained problem
\begin{equation}\label{prob:split}
  \min_{\bx,\by}\; f(\bx)+h(\by)\quad\text{s.t.}\quad \calA(\bx)=\by,
\end{equation}
which brings it within the scope of penalty methods and augmented Lagrangian methods
(ALMs). For general nonlinearly constrained nonconvex problems, both families control
feasibility through the penalty parameter. Quadratic penalty methods increase the
penalty without bound~\cite{kong2019complexity,lin2022complexity}. ALMs were
originally introduced to avoid such unbounded penalties
\cite{hestenes1969multiplier,powell1969method}, and in convex programming any fixed
penalty suffices~\cite{rockafellar1973multiplier}; in the nonconvex setting, however,
inexact ALMs still increase the penalty across outer iterations to obtain complexity
guarantees~\cite{sahin2019inexact,li2021rate}, and proximal ALM, which keeps the
penalty fixed within a run, chooses its value according to the target accuracy
\cite{xie2021proximal}. Single-loop methods have attracted growing attention because
they avoid solving subproblems to high accuracy, yet they rely on the same mechanism:
single-loop quadratic penalty schemes increase the penalty~\cite{alacaoglu2024complexity},
and linearized ALMs use small dual steps oraccuracy $eps$-dependent parameters
\cite{shi2025momentum}. A large penalty inflates the smoothness constant of the
primal update, which forces small step sizes and makes the primal subproblems
increasingly ill-conditioned.

The split formulation~\eqref{prob:split} also carries favorable structure. Its
constraint Jacobian $[J_{\calA}(\bx),\,-I]$ has full row rank at every $\bx$, so
constraint qualifications hold automatically, and its objective separates into a
smooth part in $\bx$ and a prox-friendly part in $\by$. Nonconvex ADMM and related
Lagrangian-based splitting methods exploit exactly this structure
\cite{bot2020proximal,cohen2022dynamic,hallak2023adaptive,elbourkhissi2025perturbed}.
The structure, however, does not remove the need for a large penalty. For nonlinear
$\calA$, adaptive Lagrangian schemes require penalty levels proportional to
$1/\epsilon$~\cite{hallak2023adaptive}, and linearized perturbed ALMs use
accuracy-dependent penalty and perturbation parameters~\cite{elbourkhissi2025perturbed}.
The nonconvex--linear minimax route attains $\mathcal O(\epsilon^{-3})$
\cite{pan2021efficient,he2024approximation,xu2026rada}, but its dual regularization
acts as an effective penalty that grows as the target accuracy decreases. Fixed-penalty
guarantees are available for smooth nonlinear equality constraints under a uniform LICQ
\cite{elbourkhissi2025linearized,liu2026fixed}, but they do not cover the nonsmooth
term $h(\by)$ in~\eqref{prob:split}. It therefore remains open whether a single-loop
ALM can keep every penalty-type parameter independent of both the accuracy
and the target accuracy while retaining an explicit first-order update for
$h\circ\calA$.

From another perspective, Problem~\eqref{prob:main} also admits a minimax reformulation. In particular, eliminating $\by$ from the Lagrangian of~\eqref{prob:split} gives
\begin{equation}\label{eq:ncl}
  \min_{\bx\in\RR^n}\max_{\blmbd\in\RR^m}\;
  \calL(\bx,\bq) =  \Phi(\bx,\blmbd)-h^*(\blmbd), \qquad \Phi(\bx,\blmbd):=f(\bx)+\langle\blmbd,\calA(\bx)\rangle,
\end{equation}
where $\operatorname{dom}h^*\subseteq\mathbb B(\mathbf 0,L_h)$ by
Assumption~\ref{ass:h}. This is a nonconvex--linear instance of the smooth
nonconvex--concave (NC-C) minimax problem. Since $\blmbd$ is exactly the multiplier of $\calA(\bx)=\by$, KKT points
of~\eqref{prob:split} and game-stationary (GS) points of~\eqref{eq:ncl} agree up to
constants, and ALMs can be read as single-loop minimax methods. 
For large-scale problems, we focus on single-loop first-order methods, which
update each variable by one explicit step per iteration without an inner
solver. The basic method of this kind for~\eqref{eq:ncl} is gradient
descent--ascent (GDA). It converges slowly because regularity fails on both
sides: $\Phi$ is only linear in $\blmbd$, so the best response
$\blmbd^\star(\bx)$ is in general set-valued, and $\Phi$ is nonconvex in $\bx$,
so $\bx^\star(\blmbd)$ need not be unique or stable. Even with two timescales,
GDA needs $\mathcal O(\epsilon^{-6})$ iterations to reach an $\epsilon$-GS
point~\cite{lin2020gradient}.

Two techniques improve on this rate, each restoring regularity on one side.
The first is primal smoothing, which adds $\frac{1}{2\gamma}\|\bx-\bz\|^2$ with
a fixed $\gamma$ and moves the center slowly by
$\bz^{t+1}=\bz^t+\beta(\bx^{t+1}-\bz^t)$. The target accuracy enters only
through $\beta$: with $\beta=\Theta(\epsilon^2)$, the resulting Smoothed GDA
finds an $\epsilon$-GS point in $\mathcal O(\epsilon^{-4})$ iterations. For the
pointwise maximum of finitely many smooth functions, i.e.,~\eqref{eq:ncl} with
$h(\bu)=\max_i u_i$, strict complementarity permits a constant $\beta$ and
improves the bound to $\mathcal O(\epsilon^{-2})$~\cite{zhang2020single}.

The second technique is dual regularization, which subtracts
$\frac{1}{2s}\|\blmbd-\bmu\|^2$ from the dual objective and makes it strongly
concave. In our reduced augmented-Lagrangian formulation, this regularization
has an effective-penalty interpretation:
\begin{equation}\label{eq:intro-effective-penalty}
  \max_{\blmbd}\mathcal P_s(\bx,\blmbd;\bz,\bmu)
  =\widehat{\calL}_{\rho+s}(\bx,\bmu)
    +\frac{1}{2\gamma}\|\bx-\bz\|^2,
\end{equation}
where $\mathcal P_s$ is defined in~\eqref{eq:regularized-potential},
with its dependent $\by$-argument suppressed, and
$\widehat{\calL}_{\rho+s}$ is the augmented Lagrangian minimized over
$\by$; the derivation appears in Section~\ref{sec:dual-penalty-identity}.
A zero center $\bmu=\mathbf 0$ recovers a quadratic penalty with
coefficient $\rho+s$, which can leave a nonzero feasibility residual
when $s$ is fixed. Updating $\bmu$ instead provides a multiplier-update
mechanism that may recover feasibility while keeping $\rho+s$ fixed;
its convergence and complexity depend on the coupled updates.

With the dual center fixed at $\bmu=\mathbf 0$, Perturbed GDA
attains an $\mathcal O(\epsilon^{-4})$ complexity for finding an
$\epsilon$-GS point~\cite{lu2020hybrid,xu2023unified}, and Li et al. establish
a matching algorithm-specific lower bound~\cite{li2026smoothing}.
Methods exploiting the NC-L structure attain
$\mathcal O(\epsilon^{-3})$~\cite{pan2021efficient,he2024approximation,xu2026rada}.
These guarantees rely on accuracy-dependent regularization parameters,
corresponding to an effective penalty that increases as $\epsilon$ decreases. To improve the iteration complexity, Li et al. combine primal smoothing
with zero-centered dual perturbation to develop Perturbed Smoothed GDA,
which yields a genius improvement in complexity order, reducing
the GS complexity of Perturbed GDA from $\mathcal O(\epsilon^{-4})$
to $\mathcal O(\epsilon^{-3})$ with a matching lower
bound~\cite{li2026smoothing}. However, the dual regularization coefficient still satisfies
$1/s=\Theta(\epsilon)$.

With a moving dual center $\bmu$ and fixed primal and dual regularization
coefficients, DS-GDA achieves $\mathcal O(\epsilon^{-4})$ complexity
for nonconvex--concave problems. Faster rates are available under a
uniform one-sided K{\L} condition~\cite{zheng2023universal}.
Without primal smoothing, PPALA reports
$\widetilde{\mathcal O}(\epsilon^{-2})$ complexity with fixed parameters
for functionally constrained optimization~\cite{kim2024fast}.
However, its feasibility argument relies on asymptotic reasoning,
and a corresponding finite-iteration feasibility bound remains
to be verified.  These results motivate the following question within the
primal-smoothed single-loop framework:
\blue{
\begin{quote}
  \textbf{(Q1)}
  \emph{With a moving dual center, can an $\mathcal O(\epsilon^{-3})$
  complexity for finding an $\epsilon$-KKT point be achieved with
  $s$ independent of the target accuracy, or is growth of $s$ necessary?}
\end{quote}}

In this paper, our current result answers the achievability part of (Q1) in the
increasing-$s_k$ regime. Table~\ref{tab:known-results} summarizes representative complexity guarantees, contrasting accuracy-dependent schedules with fixed base
parameters and identifying the role of primal smoothing and the dual
update mechanism.

% The fixed-$s$ question, including the role of primal smoothing and structural conditions that may yield faster rates, remains the subject of the research objectives stated below.

\begin{table*}[tp]
\centering
\begin{threeparttable}
\scriptsize
\setlength{\tabcolsep}{4pt}
\renewcommand{\arraystretch}{1.2}
\begin{tabular}{@{}lcccccc@{}}
\toprule
\multirow{3}{*}{Algorithms}
 & \multicolumn{4}{c}{Regularization and penalty parameters}
 & \multicolumn{2}{c}{Complexity} \\
\cmidrule(lr){2-5}\cmidrule(lr){6-7}
 & \multirow{2}{*}{\makecell{Penalty\\parameter}}
 & \multirow{2}{*}{\makecell{Primal smoothing\\$\frac{1}{2\gamma}\|\bx-\bz\|^2$}}
 & \multicolumn{2}{c}{\makecell{Dual regularization\\$-\frac{1}{2s}\|\blmbd-\bmu\|^2$}}
 & \multirow{2}{*}{Upper bound} & \multirow{2}{*}{Lower bound} \\
\cmidrule(lr){4-5}
 & & & Center $\bmu$ & Parameter $s$ & & \\
\midrule
\rowcolor{gray!15}
\multicolumn{7}{@{}l}{\textbf{GDA-based Methods}} \\[2pt]
\makecell[l]{\emph{Smoothed GDA}\\\cite{zhang2020smoothed,li2026smoothing}}
 & --- & $\gamma=\Theta(1)$ & --- & ---
 & $\mathcal O(\epsilon^{-4})$ & $\Omega(\epsilon^{-4})$ \\[3pt]
\makecell[l]{\emph{Perturbed GDA}\\\cite{lu2020hybrid,xu2023unified,li2026smoothing}}
 & --- & ---
 & $\bmu\equiv\mathbf 0$ & $\Theta(\epsilon^{-1})$
 & $\mathcal O(\epsilon^{-4})$ & $\Omega(\epsilon^{-4})$ \\[3pt]
\makecell[l]{\emph{Perturbed Smoothed GDA}\\\cite{li2026smoothing}}
 & --- & $\gamma=\Theta(1)$
 & $\bmu\equiv\mathbf 0$ & $\Theta(\epsilon^{-1})$
 & $\mathcal O(\epsilon^{-3})$ & $\Omega(\epsilon^{-3})$ \\[3pt]
\makecell[l]{\emph{RADA}\\\cite{xu2026rada}}
 & --- & ---
 & $\bmu\not\equiv\mathbf 0$ & $(r_\epsilon+\beta_k)^{-1}$
 & $\mathcal O(\epsilon^{-3})$ & --- \\[3pt]
\makecell[l]{\emph{DS-GDA}\\\cite{zheng2023universal}}
 & --- & $\gamma=\Theta(1)$
 & $\bmu\not\equiv\mathbf 0$ & $\Theta(1)$
 & $\underline{\mathcal O(\epsilon^{-4})}$ & --- \\
\midrule
\rowcolor{gray!15}
\multicolumn{7}{@{}l}{\textbf{Penalty and augmented Lagrangian-based methods}} \\[2pt]
\makecell[l]{\emph{Quadratic penalty}\\\cite{alacaoglu2024complexity}}
 & $\rho_k=\Theta((k+1)^{1/4})$ & ---
 & $\bmu\equiv\mathbf 0$ & $\Theta((k+1)^{1/4})$
 & $\widetilde{\mathcal O}(\epsilon^{-4})$ & --- \\[3pt]
% \makecell[l]{\emph{MLALM}\\\cite{shi2025momentum}}
%  & \makecell{$\rho=\Theta(K^{1/4})$\\$=\Theta(\epsilon^{-1})$}
%  & --- & --- & ---
%  & $\mathcal O(\epsilon^{-4})$ & --- \\[3pt]
\makecell[l]{\emph{GDPA}\\\cite{luSingleLoopGradientDescent2022}}
 & $\rho_k=\Theta(\epsilon^{-1})$
 & ---
 & $\bmu\equiv\mathbf 0$ & $\Theta(\epsilon^{-1})$
 & $\mathcal O(\epsilon^{-3})$ & --- \\[3pt]
\makecell[l]{\emph{PPALA}\\\cite{kim2024fast}}
 & $\rho=\Theta(1)$ & ---
 & $\bmu\not\equiv\mathbf 0$ & $\Theta(1)$
 & $\widetilde{\mathcal O}(\epsilon^{-2})$ & --- \\[3pt]
 \midrule
\makecell[l]{\textbf{\emph{SPLALM}}\\\emph{Theorem~?}}
 & $\rho=\Theta(1)$ & $\gamma=\Theta(1)$
 & $\bmu\not\equiv\mathbf 0$ & $\Theta(1)$
 & ? & ? \\ 
\makecell[l]{\textbf{\emph{SPLALM}}\\\emph{Theorem~\ref{thm:main}}}
 & $\rho=\Theta(1)$ & $\gamma=\Theta(1)$
 & $\bmu\not\equiv\mathbf 0$ & $\Theta(\epsilon^{-1})$
 & $\mathcal O(\epsilon^{-3})$ & ? \\
\bottomrule
\end{tabular}
\caption{Representative single-loop methods, grouped by their original
formulation. The constrained rows concern nonlinear constraints. Upper bounds display the dependence on $\epsilon$, suppressing
problem constants and bounded initial gaps. The regularization columns specify parameters and centers, not gradient
stepsizes; primal smoothing uses a separate
moving center. ``---'' denotes an absent regularizer or an unreported lower
bound; question marks denote results to be established. Constant parameters
are independent of $\epsilon$ and $k$. The rates use each method's stated
stationarity criterion and assumptions; lower bounds are algorithm-specific.}
\label{tab:known-results}
\end{threeparttable}
\end{table*}
\clearpage

The second question isolates the contribution of primal smoothing. Existing
results show a complexity improvement from primal smoothing in the
zero-centered, accuracy-dependent perturbation regime. Under a fixed
$s$, a moving dual center, and the same approximate KKT criterion, the
corresponding comparison with a method without primal smoothing is not
established.
\begin{quote}
  \textbf{(Q2)}
  \emph{Under a fixed $s$ and a moving dual center, what complexity gain,
  if any, is provided by primal smoothing?}
\end{quote}

The third question concerns structural conditions that can improve the
fixed-$s$ rate. A global error bound linking the dual-center error to the
solution of the primal proximal subproblem could reduce the amount of dual
tracking required at each iteration. Such conditions are known in special
settings, including linearly constrained problems and pointwise maxima, but
their analogue for the present composite problem with nonlinear $\calA$
remains to be characterized.
\begin{quote}
  \textbf{(Q3)}
  \emph{Which structural conditions on $(f,h,\calA)$ yield a center--primal
  solution error bound under fixed $s$, and what complexity does it imply?}
\end{quote}

Our contributions are organized as follows.
\begin{enumerate}
    \item \textbf{A single-loop method with an
    $\mathcal O(\epsilon^{-3})$ complexity guarantee.}
    We develop SPLALM by exploiting the proximal structure of the
    splitting variable $\by$. The method combines primal smoothing,
    centered dual regularization, and moving primal and dual centers,
    without an iterative inner solver.
    With $s_k=\Theta((k+1)^{1/3})$ and suitable center stepsizes,
    we establish an $\mathcal O(\epsilon^{-3})$ iteration bound
    for finding an $\epsilon$-KKT point.
    The base penalty $\rho$ and the primal and dual gradient stepsizes
    remain fixed, while the effective penalty $\rho+s_k$ increases.

    % 已有：Algorithm 3.1、参数调度 (3.9)、Theorem 3.1。
    % 已有：第 2 节的正则化--有效罚参数恒等式与零中心偏差例子。
    % 零中心例子证明：在该实例上，正则化问题的精确驻点及其
    % 近端证书达到 epsilon-KKT，需要 rho+s = Omega(epsilon^{-1})。
    % 这属于驻点偏差结论，不属于迭代复杂度下界。
    %
    % TODO：统一算法名称。当前 PDF 使用 SPLALM，
    % 旧研究主线使用 DDLALM。
    %
    % TODO：研究 Q2，即 primal smoothing 的独立贡献。
    % 首先核实 PPALA 的具体版本、假设和有限步可行性证明。
    % 在相同问题假设、固定 s、移动 mu 和相同输出准则下，
    % 分析去掉 primal smoothing 的版本。
    % 固定 s 与递增 s_k 分别比较，不交叉归因。
    % 区分“改善已知上界”与“证明实际最坏复杂度的分离”。
    % 该部分目前不是现有定理的直接推论。

    \item \textbf{The complexity of ordinary center updates
    under a fixed effective penalty.}

    % TODO：首先建立固定 s 版本的复杂度基准。
    % 核实 DS-GDA 分析向当前 D_rho 模型迁移所需的假设，
    % 包括迭代有界性、光滑常数和更新顺序。
    % 目标上界：O(epsilon^{-4})；尚未证明前不列为本文结果。
    %
    % TODO：证明实际单环更新下的乘子有界性。
    % 利用 q^k in partial h(y^{k+1})、||q^k|| <= L_h，
    % 以及 tau <= rho*s/(rho+s) 下的凸组合结构。
    % 原始迭代有界性需另证，不能仅从目标水平集有界直接推出。
    %
    % TODO：研究固定 s 上界是否紧；下界正文暂留为空。
    % 明确下界覆盖的更新族、参数调度、oracle 和输出规则，
    % 保留初始间隙 Delta 及问题常数的依赖。
    % Li et al. 的困难实例仅作为候选起点；
    % 必须核实它能否嵌入当前复合问题及实际更新公式。
    %
    % TODO：若建立匹配的 Omega(epsilon^{-4}) 下界，
    % 则与递增 s_k 的 O(epsilon^{-3}) 上界形成分离。
    % 结论只限于下界覆盖的更新族和一般问题类。
    % 不写成“所有 ALM 都必须增大罚参数”，
    % 也不把 s 增长称为“唯一必要条件”。
    %
    % TODO：数值实验分别比较：
    % 相同 s 调度下，零中心与移动中心的效果；
    % 相同中心规则下，固定 s 与递增 s_k 的效果。
    % 同时报告驻点残差、可行性残差和实际 oracle 调用次数。
    %
    % 可选：研究加速中心更新。
    % 对固定原始中心的精确子问题，加速近端点法提供启发；
    % 对当前单环算法，需重新控制原始中心漂移和求解误差。
    % 暂不声称加速后必然达到 O(epsilon^{-3})。

    \item \textbf{Structural conditions for faster convergence
    with a fixed effective penalty.}

    % TODO：定义原始近端子问题及其中心--原始解误差界。
    % 明确 X_s(z,mu)、X_*(z)、Lambda_s(z,mu) 的存在唯一性，
    % 以及误差界常数在哪个集合上一致成立。
    %
    % 候选误差界：
    % ||X_s(z,mu)-X_*(z)||^2
    %     <= C ||Lambda_s(z,mu)-mu||^p.
    %
    % TODO：证明该误差界如何控制实际单环迭代的跟踪误差，
    % 并通过势函数与残差传递得到 epsilon-KKT 复杂度。
    % O(epsilon^{-4/p}) 暂为待验证目标；
    % p=2 的端点需单独分析，目标为 O(epsilon^{-2})。
    %
    % TODO：研究可验证的充分条件：
    % 线性 A 与适当的秩或误差界条件；
    % 严格互补与有效方向非退化的组合条件；
    % 原始近端 KKT 系统的适当度量次正则性。
    %
    % 对非线性 A，不能仅凭 J_A J_A^T 的一致正定性
    % 就断言近端对偶一致强凹；需要证明相应的统一曲率或误差界。
    % primal smoothing 保证原始子问题强凸时，
    % 也不自动保证去掉对偶正则项后的近端对偶强凹。
    %
    % TODO：区分局部与全局结果。
    % 若条件仅局部成立，应证明迭代进入并保持在该邻域，
    % 或明确给出局部初始化条件。
    %
    % 可选扩展：去掉 h 的 Lipschitz 假设。
    % 必须以其他条件重新建立乘子有界性与所需光滑性，
    % 不直接从 Jacobian 非退化推断全过程乘子有界。
\end{enumerate}

\begin{table*}[t]
\centering
\small
\caption{Current results and research objectives. Here
$\rho+s_k$ is the effective penalty.
The bounds in the last two rows are research targets.}
\label{tab:research-roadmap}
\setlength{\tabcolsep}{5pt}
\renewcommand{\arraystretch}{1.15}
\begin{tabular}{@{}p{0.19\textwidth}p{0.22\textwidth}
                  p{0.28\textwidth}p{0.22\textwidth}@{}}
\toprule
Setting & Center and penalty & Result or objective & Status \\
\midrule
Zero-centered formulation
& $\bmu=\mathbf 0$, $\rho+s$
& A stationary-point example requires
  $\rho+s=\Omega(\epsilon^{-1})$
& Feasibility-bias example in Section~2 \\
\addlinespace
Current SPLALM
& Moving $\bmu$;
  $s_k=\Theta((k+1)^{1/3})$
& $\mathcal O(\epsilon^{-3})$ for $\epsilon$-KKT
& Theorem~\ref{thm:main} \\
\addlinespace
Fixed-$s$ variant
& Moving $\bmu$; fixed $\rho+s$
& Establish an $\mathcal O(\epsilon^{-4})$ bound
  and investigate its tightness
& Pending analysis \\
\addlinespace
Fixed-$s$ variant with structure
& Moving $\bmu$; fixed $\rho+s$
& Derive faster rates from a proximal error bound;
  target $\mathcal O(\epsilon^{-2})$
& Pending analysis \\
\bottomrule
\end{tabular}
\end{table*}


Structure of the paper. The paper is organized as follows. In Section 2, we introduce the preliminaries,
including assumptions, stationarity notions, and the two acceleration tools: dual perturbation and smoothing.
The main body of the paper focuses on single-loop methods: In Section 3, we present the main results,
including the proposed Perturbed Smoothed GDA together with all iteration-complexity bounds and their
tight analyses. In Section 4, we develop a unified convergence analysis framework for these single-loop
algorithms. In Section 5, we provide a tight analysis that sharpens these iteration complexities. In Section 6,
wefurthershowthatthetwo accelerationtechniquesand theunifiedanalysisfrom Section4also applyto a
double-loop setting, resulting in a new algorithm, Perturbed Smoothed FOAM, improving the best known
iterationcomplexities. Finally,Section7concludes the paper